%MIT OpenCourseWare: https://ocw.mit.edu
%RES.18-012 Algebra II Student Notes, Spring 2022
%License: Creative Commons BY-NC-SA 
%For information about citing these materials or our Terms of Use, visit: https://ocw.mit.edu/terms.

% \rhead{\rightmark}
\lhead{Lecture 10: Ideals in Polynomial Rings}

\section{Ideals in Polynomial Rings}

\subsection{Ideals in a Field}

Last class, we looked at ideals in rings, and began discussing ideals in polynomial rings. Today we will consider what such ideals can look like. The following proposition will be useful later, when discussing maximal ideals:

\begin{proposition} \label{field iff two ideals}
    A ring $R$ is a field if and only if it has exactly two ideals. 
\end{proposition}

Any ring has the ideals $(0) = \{0\}$ and $(1) = R$ --- these coincide only in the zero ring, which is not a field by definition. So the proposition states that the ring is a field if and only if it has no other ideals. 

\begin{proof}
    Suppose $R$ is a field, so it has at least two ideals $(0)$ and $(1)$. But there are no other ideals, because every element is invertible --- if $I$ is an ideal containing some element $x \neq 0$, then $1 = x^{-1}x$ is in $I$ as well, so $I = (1)$. 

    Conversely, if $R$ is not a field, then either it is the zero ring and only has one ideal, or it contains a nonzero $x$ which is not invertible. Then $(x)$ cannot contain $1$, so $(0)$, $(x)$, and $(1)$ are distinct ideals. 
\end{proof}

\subsection{Polynomial Rings over a Field}

We'll first look at ideals in $F[x]$, the ring of polynomials in one variable over a field. 

\begin{proposition} \label{F[x] is a PID}
    Every ideal in $F[x]$ is principal. More precisely, if $I \subset F[x]$ is a nonzero ideal and $P$ a (nonzero) element of $I$ of minimal degree, with $\deg(P) = n$, then we have $I = (P)$, and the images of $1$, $x$, $x^2$, \ldots, $x^{n - 1}$ form a basis in $F[x]/I$ (as a vector space over $F$). 
\end{proposition}

\begin{proof}
    The main idea is to use division of polynomials with remainder. 
    
    In order to check that $P$ generates $I$, take any $Q \in I$; then by polynomial division we can write \[Q = P\cdot S + R,\] where $S$ and $R$ are in $F[x]$ and $\deg R < \deg P$. But $R$ must be in $I$, so if $R \neq 0$ then this contradicts the choice of $P$ as having minimal degree. So $R = 0$, which means $P \mid Q$. So $P$ generates $I$. 
    
    Now consider the quotient $F[x]/(P)$, and let the images of $1$, $x$, \ldots, $x^{n - 1}$ be denoted by $\overline{1}$, $\overline{x}$, \ldots, $\overline{x^{n - 1}}$. These images must be linearly independent --- for any $P, Q \in F[x]$ we have $\deg(PQ) = \deg(P) + \deg(Q)$ (since $F$ is a field, so has no zero divisors), which means a polynomial with degree less than $n$ cannot be divisible by $P$. Then if $\overline{1}$, \ldots, $\overline{x^{n - 1}}$ were linearly dependent with $a_0\cdot \overline{1} + \cdots + a_{n - 1}\cdot \overline{x^{n - 1}} = 0$, then $a_0 + \cdots + a_{n - 1}x^{n - 1}$ would have to be divisible by $P$, which is a contradiction. 
    
    Conversely, these images must span the quotient by using division with remainder again --- for any $Q \in F[x]$, we can write $Q = P \cdot S + R$ with $\deg(R) < n$, which means $\overline{Q} = \overline{R}$ for some $\overline{R}$ which is a linear combination of $\overline{1}$, \ldots, $\overline{x^{n - 1}}$.
\end{proof}

Recall that to divide a polynomial $Q$ by $P$ with remainder, we subtract a multiple of $P$ from $Q$ to cancel out the leading term of $Q$; we then repeat until the remaining polynomial has degree less than that of $P$. 

\begin{note}
    This doesn't generalize fully to polynomials over an arbitrary ring $R$, but some parts do. First, it's not always true that \[\deg PQ = \deg P + \deg Q.\] For example, in $\mathbb{Z}/4[x]$, $(2x)(2x + 1) = 2x$ does not have degree $2$.
    
    Division with remainder also does not necessarily work --- for example, we can't divide $x^2$ by $2x + 1$ with remainder in $\mathbb{Z}[x]$. This is because when we cancel out the leading coefficient of $Q$, we need to scale; and here, $2$ isn't invertible, so we can't scale by the correct factor.

    But both facts remain true for monic polynomials --- so we can divide with remainder if $P$ is monic (meaning it has leading coefficient $1$; this works the same way if the leading coefficient is a unit). So if $P$ is monic, then it's still true that every element of $R[x]/(P)$ can be written uniquely as \[a_0\overline{1} + a_1\overline{x} + \cdots + a_{n - 1}\overline{x^{n - 1}},\] for $a_i \in R$. We no longer say that the $\overline{x^i}$ form a basis, since $R[x]/(P)$ is not a vector space (vector spaces are defined over a field); but we will later discuss an analog of vector spaces over rings.
\end{note}

Last time, we mentioned that the construction $F[x]/(P)$ is equivalent to adjoining a root of $P$ --- for example, $\CC \cong \RR[x]/(x^2 + 1)$. Using Proposition \ref{F[x] is a PID}, we can now make this more precise: 

\begin{proposition}
If $F$ is a field and $P \in F[x]$, then $F[x]/(P) \cong F[\alpha]$, where $\alpha$ is a root of $P$. 
\end{proposition}

\begin{proof}
We know that \[F[x]/(P) = \left\{a_0 + a_1\overline{x} + \cdots + a_{n - 1}\overline{x^{n - 1}}\right\},\] where $n = \deg(P)$, while \[F[\alpha] = \left\{a_0 + a_1\alpha + \cdots + a_{n - 1}\alpha^{n - 1}\right\}\] (note that we don't need higher powers of $\alpha$, as the polynomial relation guarantees that we can express them in terms of these powers). To multiply in $F[\alpha]$, it's enough to understand how to multiply by $\alpha$. We have \[\alpha(a_0 + \cdots + a_{n - 1}\alpha^{n - 1}) = a_0\alpha + \cdots + a_{n - 2}\alpha^{n - 1} + a_{n - 1}\alpha^n,\] and we can expand $\alpha^n$ using the fact that $P(\alpha) = 0$ --- if $P(x) = c_nx^n + \cdots + c_0$, then we have \[\alpha^n = -c_n^{-1}(c_{n - 1}\alpha^{n - 1} + \cdots + c_0).\] This is exactly how multiplication works in $F[x]/(P)$; so the two rings have the same additive and multiplicative structure, and are therefore isomorphic, with the isomorphism given by replacing $\overline{x}$ with $\alpha$. 
\end{proof}

\begin{note}
Here we should think of $\alpha$ as a \emph{universal} root of $P$ --- an element that satisfies $P(\alpha) = 0$ but no extra conditions. Given a \emph{specific} root, it's possible that $\alpha$ satisfies additional relations (if $P$ is not irreducible), in which case $F[\alpha]$ would instead be isomorphic to a quotient of $F[x]/(P)$. 
\end{note}

% \todo{add the reveiew section}
% \subsection{More About Ideals}
% Now, let's consider ideals in $F[x]$ where $F$ is a field. 
% \begin{proposition}
% \begin{enumerate}[label = (\alph*)]
%     \item Every ideal in $F[x]$ is  principal
%     \item If $I \subset F[x]$ is a nonzero ideal and $P \in I$ is a nonzero element of smallest degree $n$, then $I = (P)$ ad the images of $1, x, \cdots, x^{n-1}$ form a basis in $F[x]/I.$
% \end{enumerate}
% \end{proposition}
% \begin{proof}
% First, it is possible to show that part b) implies part a). Let $P \in I$ be an element of smallest degree $n.$ Then, we must check that $P$ generates $I.$ For some polynomial $Q \in I$, then by polynomial division with remainder, $Q = P\cdot S + R$ for $S, R \in F[x].$ However, since $P$ is the nonzero element of $I$ with smallest degree, the fact that $\deg(R) < \deg(P)$ implies that $R = 0.$ \todo{explain more}. Also, $\overline{1}, \overline{x}, \cdots, \overline{x}^{n+1}$ are linearly independent. The degree $\deg(PQ) = \deg(P) + \deg(Q),$ so a polynomial of degree less than $n$ is \emph{not} divisible by $P.$ \todo{why is this helpful lol} Moreover, they span the quotient because $Q = P \cdot S + R$ implies that $\overline{Q} = \overline{R}$ since $\deg(R) < n.$
% \end{proof}

% \begin{qq}\todo{make the qqs better}
% Does this generalize to $R[x]$ for an arbitrary $R$?
% \end{qq}
% No. In general, it is not true that $\deg(PQ) = \deg(P) + \deg(Q),$ and division by remainder also does not necessarily exist. For example, in $\ZZ_4[x]$, $(2x) \cdot (2x+1) = 2x$, and in $\ZZ[x],$ it is not possible to divide $x^2$ by $2x + 1$ with remainder. However, it \emph{is} possible to divide with remainder by a monic polynomial with leading coefficient 1. 

% \begin{proposition}
% If $P$ is monic, then every element in $R[x]/(P)$ can be uniquely written as $a_0\overline{1} + a_1\overline{x} + \cdots + a_{n-1}\overline{x}^{n-1}$ for $a_i \in R.$
% \end{proposition}

% Now, it is finally possible to prove this proposition from last time.
% \begin{proposition}\todo{verify that this is actually what i watn LOL}
% For a field $F,$ $F[x]/(P)$ is equivalent to adjoining a root of $P.$ \todo{ask for clarification on this}
% \end{proposition}

% For example, the complex numbers can be constructed by $\CC = \RR[x]/(x^2+1).$ 

% \begin{proof}
% Consider $R[x]/P$ and $\alpha = \overline{x}.$\todo{what does this mean} We get $\{a_0 + a_1\alpha + \cdots + a_{n-1}\alpha^{n-1}\}$ added componentwise to multiply by $\alpha.$ Writing it as $a_0\alpha + \cdots + a_{n-2}\alpha^{n-1} + a_{n-1}\alpha^n$\footnote{Euler called these structures "rings" because of the type of circle structure demonstrated here.} and expanding it using $P(d) = 0$ provides our result. \todo{why does it do that LOLLLLL}
% \end{proof}

\subsection{Maximal Ideals}

We'll now discuss maximal ideals. These will turn out to be quite useful; today we'll see how to use them to connect algebra to geometry. 


% We now discuss a particular type of ideal.\todo{make this sound less stupid}

\begin{definition}[Maximal Ideals]
An ideal $I \subset R$ is \textbf{maximal} if $I \neq R$, and the only ideals of $R$ containing $I$ are $R$ and $I$ itself. 
\end{definition}

\begin{example}
    The maximal ideals in $\ZZ$ are $(p)$, for $p$ prime. To prove this, we saw earlier that the ideals of $\mathbb{Z}$ are $n\mathbb{Z}$. But $n\mathbb{Z} \subset m\mathbb{Z}$ if and only if $m \mid n$. So understanding the maximal \emph{ideals} of $\ZZ$ in the poset of ideals ordered by inclusion is equivalent to understanding the minimal \emph{elements} of $\ZZ$ in the poset of elements ordered by divisibility. These minimal elements are the primes $p$, so the maximal ideals are $(p)$. 
\end{example}
% \begin{example}
% What are the maximal ideals in $\ZZ$? Clearly, the statement that $n\ZZ \subset m\ZZ$ implies that $m$ divides $n$, and so understanding the maximal \emph{ideals} of $\ZZ$ is equivalent to understanding the minimal dividing \emph{elements} in $\ZZ.$ In particular, the maximal ideals in $\ZZ$ are $(p)$ for $p$ prime.
% \end{example}

\begin{example}
For a polynomial ring over a field $F[x]$, any ideal is of the form $(P)$. For the same reason as in the case of $\ZZ$, the ideal $(P)$ is maximal if and only if $P$ does not factor as $QR$ where $Q$ and $R$ have positive degree, or in other words, if $P$ is irreducible. 
\end{example}

\begin{example}
In $\CC[x]$, the only irreducible polynomials are linear, since by the Main Theorem of Algebra every polynomial can be factored as $c(x - z_1)\cdots (x - z_n)$. So the maximal ideals of $\CC[x]$ are exactly the ideals $(x - \alpha)$. 
\end{example}

We'll now see how this example generalizes to polynomials in \emph{multiple} variables. The following proposition will be useful:

% Now, we see how this generalizes.\todo{i missed what he said}
\begin{proposition} \label{maximal iff field}
An ideal $I \subset R$ is maximal if and only if $R/I$ is a field.
\end{proposition}

\begin{proof}
First, $R/I$ is a field if and only if $R/I$ has exactly two ideals (by Proposition \ref{field iff two ideals}). But by the correspondence theorem for rings, ideals in $R/I$ are in bijection with ideals in $R$ containing $I$. So $R/I$ is a field if and only if $R$ has exactly two ideals containing $I$. But this is equivalent to the condition that $I$ is maximal (since $I$ and $R$ are both containing $I$, so if there are only two such ideals, then there can be no others).
\end{proof}

\begin{example}
In the case of $\ZZ$, as we've mentioned earlier, $\ZZ/p\ZZ$ is a field.
\end{example}

\begin{example}
If $F$ is a field and $P \in F[x]$ is irreducible, then $F[x]/(P)$ is a field as well --- this is a construction which can be used to build new fields. 
\end{example}

This describes what happens when we look at polynomials in \emph{one} variable over a field, but we can try to consider what happens for polynomials in \emph{multiple} variables as well. 

\subsection{Ideals in Multivariate Polynomial Rings}

% \subsection{Examples}\todo{make all the sections/subsections better titled LOL}
% This is a central idea in ring theory. Let's see some examples!

\vspace{0.75em}

\begin{example}
    Let $R = F[x_1, \ldots, x_n]$, where $F$ is a field. Fixing scalars $\alpha = (\alpha_1, \ldots, \alpha_n)$ (with $\alpha_i \in F$), we have the evaluation homomorphism $F[x_1, \ldots, x_n] \to F$ defined as \[\operatorname{ev}_\alpha : P \mapsto P(\alpha_1, \ldots, \alpha_n).\] This map is clearly onto, as the constants in $F[x_1, \ldots, x_n]$ are sent to themselves, so by the first isomorphism theorem for rings, we have \[F \cong F[x_1, \ldots, x_n]/\ker(\operatorname{ev}_\alpha).\] Since $F$ is a field, the kernel is a maximal ideal of $R$. In fact, this kernel can be explicitly written as \[\mathfrak{m}_\alpha = (x_1 - \alpha_1, \ldots, x_n - \alpha_n).\]
\end{example}

This provides a way to construct maximal ideals of $F[x_1, \ldots, x_n]$; building on this, we can also try to construct maximal ideals of \emph{quotients} of this ring (since many rings can be constructed in this way). 

Suppose that $R = F[x_1, \ldots, x_n]/J$, with $J = (P_1, \ldots, P_m)$. (We'll later see that for \emph{any} $J$, we can find finitely many polynomials $P_i$ which generate $J$; but we won't focus on that right now.) Then for any $\alpha$ for which $\mathfrak{m}_\alpha \supset J$, where $\mathfrak{m}_\alpha$ is the maximal ideal of $F[x_1, \ldots, x_n]$ as defined above, by the correspondence theorem the image of $\mathfrak{m}_\alpha$ is also a maximal ideal of $F[x_1, \ldots, x_n]/J$. 

But we know $\mathfrak{m}_\alpha \supset J$ if and only if $\alpha = (\alpha_1, \ldots, \alpha_n)$ is a common zero of $P_1$, \ldots, $P_n$ (since $\mathfrak{m}_\alpha = (x_1 - \alpha_1, \ldots, x_n - \alpha_n)$ contains $P_i$ if and only if $\alpha$ is a zero of $P_i$). This gives the following construction:

\begin{proposition}
Suppose $R = F[x_1, \ldots, x_n]/(P_1, \ldots, P_m)$. Then any common zero $\alpha = (\alpha_1, \ldots, \alpha_n)$ of $P_1$, \ldots, $P_m$ yields a maximal ideal of $R$, which is the image of $\mathfrak{m}_\alpha$ when quotienting out by $(P_1, \ldots, P_m)$. 
\end{proposition}

The sets of common zeroes of a list of polynomials are studied in algebraic geometry. Here, we can see that such common zeroes can be used to produce maximal ideals of $R$. A famous theorem states that in the case $F = \CC$, the converse is true as well:

% \begin{example}
% Let $R = F[x_1, \cdots, x_n]$ where $F$ is a field. Fixing $\overline{\alpha} = (\alpha_1, \cdots, \alpha_n) \in F,$ the evaluation map is 
% \begin{align*}
% \text{ev}_{\overline{\alpha}}: F[x_1, \cdots, x_n] &\rto F\\
% P &\mto P(\alpha_1, \cdots, \alpha_n)
% \end{align*}

% By the first isomorphism theorem, it can be identified with the quotient, \todo{explain this}so its kernel is a maximal ideal $M_{\overline{\alpha}} = (x_1 - \alpha_1, \cdots, x_n - \alpha_n).$
% \end{example}

% Building on this example, it is also possible to take $R = F[x_1, \cdots, x_n]/J$ where $J = (P_1, \cdots, P_m)$ is generated by $P_i.$ This is pretty much a universal example, since if the ring can be found inside a field, it can be written as a quotient. \todo{what is he saying} If $M_\alpha \supset J,$ then by the correspondence theorem, there is a maximal ideal in the quotient $R$. Equivalently $\overline{\alpha} = (\alpha_1, \cdots, \alpha_n)$ is a common zero of $P_1, \cdots, P_m.$

% \todo{he had some explanation here}

% The following is a famous theorem. 
\begin{theorem}[Hilbert's Nullstelensatz]\label{nullstelensatz}
Every maximal ideal of $\CC[x_1, \cdots, x_n]$ is of the form $\mathfrak{m}_{\alpha}$ for some $\alpha = (\alpha_1, \cdots, \alpha_n)$.
\end{theorem}

In the name of the theorem (in German), ``null'' means zero, ``stelen'' means place, and ``satz'' means theorem.

This immediately implies the following corollary (since maximal ideals in $\CC[x_1, \cdots, x_n]/(P_1, \cdots, P_m)$ are in correspondence with maximal ideals in $\CC[x_1, \ldots, x_n]$ containing all the $P_i$):

\begin{corollary}
The maximal ideals in $R = \CC[x_1, \cdots, x_n]/(P_1, \cdots, P_m)$ are in bijection with the common zeroes of the polynomials $P_i$.
\end{corollary}

\begin{proof}[Proof of Theorem \ref{nullstelensatz}]
Let $\mathfrak{m} \subset \mathbb{C}[x_1, \ldots, x_n]$ be a maximal ideal; then $F = \mathbb{C}[x_1, \ldots, x_n]/\mathfrak{m}$ is a field. This gives a homomorphism from $\mathbb{C}$ to $F$ --- taking the quotient by $\mathfrak{m}$ gives a homomorphism from $\mathbb{C}[x_1, \ldots, x_n]$ to $F$, and we can restrict the homomorphism to $\CC$. It suffices to show that this map $\CC \to F$ is an isomorphism --- then the original homomorphism $\CC[x_1, \ldots, x_n] \to F$ from taking the quotient must map $\CC$ isomorphically to $F$, and we can suppose it maps $x_i \to \alpha_i$ for each $i$ where $\alpha_i \in \CC$ (it really maps $x_i$ to some element of $F$, but $F$ is isomorphic to $\CC$). Then it's clear that the kernel of this homomorphism is generated by $x_1 - \alpha_1$, \ldots, $x_n - \alpha_n$; therefore this kernel is $\mathfrak{m}_\alpha$ where $\alpha = (\alpha_1, \ldots, \alpha_n)$, and we have $\mathfrak{m} = \mathfrak{m}_\alpha$. So now we want to show that the map $\CC \to F$ is bijective. 

But \emph{any} homomorphism between fields is injective --- the kernel of the homomorphism must be an ideal, but the only ideals of a field are $\{0\}$ and the entire field. Since the homomorphism cannot map $1$ to $0$, the kernel cannot be the entire field, so must be $\{0\}$. 

So it suffices to show that the homomorphism is surjective. Assume not. Then $F$ strictly contains $\mathbb{C}$ (since we have an injective map $\CC \hookrightarrow F$, so $\CC$ is isomorphic to its image), so we can pick $z \in F$ with $z \not\in \mathbb{C}$. Then consider the elements \[\left\{ \frac{1}{z - \lambda} \mid \lambda \in \mathbb{C}\right\}.\] 

Now we'll use a bit of set theory --- $\mathbb{C}[x_1, \ldots, x_n]$ is a countable union of finite-dimensional vector spaces $U_1 \subset U_2 \subset \cdots$ over $\mathbb{C}$, where $U_i$ is the vector space of polynomials with degree at most $i$. Then since $F$ is a quotient of $\CC[x_1, \ldots, x_n]$, it must also be a countable union of finite-dimensional vector spaces $V_1 \subset V_2 \subset \cdots$, where $V_i$ is simply the image of $U_i$ in this quotient. 

On the other hand, $\mathbb{C}$ is not countable, so there are uncountably many terms $1/(z - \lambda)$. Since all such terms are elements of $F$, one of the finite-dimensional vector spaces that $F$ consists of must contain infinitely many of them. 

But then since these terms all belong to the same finite-dimensional vector space, and there's infinitely many of them, we can find a finite set which is linearly dependent --- so then we have an identity \[\sum \frac{a_i}{z - \lambda_i} = 0,\] where $a_i \in \mathbb{C}$ for all $i$ and there are finitely many terms. But by clearing denominators, we can translate this into a polynomial relation in $z$ --- we then have $P(z) = 0$ for some $P \in \CC[x]$. Then since all polynomials over $\mathbb{C}$ factor, we can write \[P(x) = c\prod_i (x - r_i),\] for some $r_i \in \mathbb{C}$. But $z \not\in \mathbb{C}$, so $z$ cannot equal any of the $r_i$, contradiction (as $F$ is a field, so the product of nonzero terms cannot be zero). 
    
So then the map $\CC \to F$ must be an isomorphism, as desired. 
\end{proof}
% Suppose $M \subset \CC[x_1, \cdots, x_n]$ is a maximal ideal. Then there is some field $F = \CC[x_1, \cdots, x_n].$ Every mapping between fields is injective. \todo{say why it is injective....} First, $\CC \hookrightarrow F$. If $\CC$ maps isomorphically to $F,$ then we are done, since $x_i \mto \alpha_i$ and then $M = (x_i - \alpha_i)$.

% On the other hand, suppose $F \supsetneq \CC$; pick $z \in F$ such that $z \notin \CC.$ Consider $\frac{1}{z - \lambda}$ for $\lambda \in \CC.$ Then $\CC[x_1, \cdots, x_n]$ is a countable union of finite dimensional space. Then $\CC$ is not a finite union of finite subsets, implying that infinitely many \todo{wtf did he put here} $\frac{1}{z-\lambda}$ fall in a function space, implying that there exists a union dependence $\sum \frac{a_i}{z - \lambda_i} = 0$ for $a_i \in \CC.$ This implies that $P(z) = 0$ for $P \in \CC[x]$, $P(x) = \prod(x - \lambda_i),$ implying that $z = \lambda_i.$ \todo{pls get this proof straightened out}
% \end{proof}